\subsection{问题三：三类患者联合排程}

\subsubsection{共享压力场景的构造}

门诊和体检数据没有预约时刻。为回答“在三类患者同时占用设备时，住院48小时率能达到多少”，本文把两类患者的历史报告日期作为回顾性计划日，构造必须在该日工作块内完成的背景任务。该场景使用患者项目组合和日期负荷，但剔除报告时分、历史机房与历史医生。因此，它反映的是“当日需求压力下的可排程度”，不是对历史派单的重放。

背景事件的基础释放时刻为历史开单时刻与计划日工作块起点的较晚者，截止为计划日工作块终点；空腹任务只能进入上午并在10时前结束。对憋尿任务，准备完成时刻取开单时刻加准备提前量，最终释放为该时刻与计划日工作块起点的较晚者；这允许已提前准备的门诊/体检患者在开班时进入排程，而不会一律延迟到9时。住院事件仍使用精确开单时刻、准备约束和48小时截止。连续日历覆盖2024年3月30日至2025年4月2日：边界前2日承接尚未逾期的在制任务，边界后2日只用于完成3月31日开单患者的48小时观察窗；评价仅统计2024年4月1日至2025年3月31日释放或计划的事件。

\subsubsection{背景服务下限与住院主目标}

题面要求在满足门诊、体检需求的前提下最大化住院48小时率，但未给出背景服务的数值阈值。若直接把背景完成数放在住院指标之前，会在背景多排少量患者时牺牲题面唯一考核指标；若完全忽略背景，又不符合业务前提。因此本文以共享FCFS得到的背景计划日完成数$B_0$作为“不比基线差”的可审计下限：
\begin{equation}
  Z_B=\sum_{e\in\Ecal_O\cup\Ecal_P}y_e\geq B_0.
\end{equation}
在满足该下限的候选策略中，首先最大化住院48小时完成数
\begin{equation}
  \max Z_I=\sum_{e\in\Ecal_I}y_e.
\end{equation}
随后才比较背景增量、住院等待、跨室和设备负荷差。共享FCFS和联合稀缺度均满足$Z_B\geq B_0$，前者住院指标更高，故作为主推荐；后者背景服务率更高，作为医院提高背景下限时的帕累托备选。若医院能给出正式门诊/体检服务水平或急诊等级，应以该阈值和急诊硬保护替代$B_0$，不能长期把历史基线当成临床制度。

主策略先按计划日与释放时刻安排背景患者，再按住院释放时刻（由于截止均为开单后48小时，亦等价于最早截止顺序）进入共享日历；每个患者内部仍采用能力稀缺优先、同室优先和最早可行时隙搜索。因此“FCFS”只描述患者间主序，不等于忽略设备能力的简单排队。联合稀缺度备选则先按空腹标记、项目最小可用设备数和工作量安排背景任务，再按住院松弛度排入余量。两种过程都只初始化一次房间占用和医生负荷状态。

\subsubsection{联合排程结果}

主情景中，主推荐策略在留出期完成57\,341个背景事件，计划日完成率59.70\%；其中门诊完成42\,871个，体检完成14\,470个。住院完成45\,260个，48小时率63.95\%。联合稀缺度备选完成59\,175个背景事件，背景率61.61\%，但住院只完成44\,805个、住院率63.31\%。两个比例的分母均包含未分配事件，未因无可行设备或搜索失败而删除。

图\ref{fig:daily}给出14日滚动完成率；14日只用于压低日级随机波动、辨认趋势，不进入排程规则。住院率大部分时间在60\%--75\%之间，门诊/体检在2024年秋季出现明显下降、2025年初回升，与体检团检波动和背景总负荷一致。年度平均值不能替代日级监控；实际部署时应同时查看完成率、未来48小时已知需求和稀缺项目分钟数，再决定是否调班。

\begin{figure}[H]
  \centering
  \includegraphics[width=\textwidth]{fig06_daily_performance.pdf}
  \caption{主推荐策略的日级服务稳定性}
  \label{fig:daily}
\end{figure}

图\ref{fig:gantt}展示主推荐策略留出期工作量最高的2024年10月14日。排程在12--13时保持午休断点，三类患者在各机房交错安排；床边B超主要服务住院患者，产科和门诊机房在上午承接体检与门诊，下午逐步释放给住院任务。图中每个色块均来自一条实际患者项目记录，不是总容量示意。

\begin{figure}[H]
  \centering
  \includegraphics[width=\textwidth]{fig07_representative_gantt.pdf}
  \caption{高负荷代表日的22台设备患者级排程}
  \label{fig:gantt}
\end{figure}

设备日均占用率见图\ref{fig:utilization}. 7号床边B超达到73.95\%，显著高于其他设备的36.16\%--57.05\%，与表\ref{tab:capability}的唯一床旁能力一致。其余普通设备存在一定负荷差异；产科III/IV级设备占用相对较低并不表示可以撤销，因为它们承担的是不可替代的专项能力。

\begin{figure}[H]
  \centering
  \includegraphics[width=0.88\textwidth]{fig08_room_utilization.pdf}
  \caption{主推荐策略下各设备日均占用率}
  \label{fig:utilization}
\end{figure}

\subsubsection{联合方案的实施建议}

由结果得到四项实施建议。

\begin{enumerate}
  \item 每次发布下一轮计划前，用未来48小时已知门诊、体检计划和住院在制队列重新构造滚动日历；已确认预约作为冻结层，新增住院患者只在余量中插入。
  \item 7号床边B超按未来48小时床旁任务上界分钟保护容量；当剩余容量已不能覆盖床旁需求或最紧迫床旁事件时，停止非床旁回填，并由医院决定是否启用具备床旁能力的备用设备或移动班次。
  \item 对64、43、65号产科高等级设备和45、44号儿科设备保留能力池，不以短期占用率为唯一裁撤依据；灵活项目优先落到能力覆盖广的普通设备。
  \item 若预计医生并发达到$c_{wt}$，优先调整医生班次而不是继续开放空机房。新增设备只有与相应医生资质和班次同时增加，才能转化为完成率。
\end{enumerate}

\section{模型的检验}

\subsection{独立可行性核验}

为避免调度器用自身逻辑证明自身正确，另写独立核验程序，只读取统一输入、最终排程和结果表，重新计算任务时长、占用区间和医生并发。核验对象为151\,020条项目任务、102\,986个已排事件。检查组见表\ref{tab:verification}，所有违规数均为0；最大观测医生并发为20，等于主容量最大值而未超出。

\begin{table}[H]
  \centering
  \caption{最终联合排程的独立核验结果}
  \label{tab:verification}
  \begin{tabularx}{\textwidth}{X r l}
\toprule
核验项目 & 违规数 & 结论 \\
\midrule
任务键、输入对应、时长 & 0 & 通过 \\
设备兼容、床旁位置 & 0 & 通过 \\
释放、截止、工作时段 & 0 & 通过 \\
房间、患者、转运冲突 & 0 & 通过 \\
医生并发容量 & 0 & 通过 \\
结果表聚合一致性 & 0 & 通过 \\
\bottomrule
\end{tabularx}

\end{table}

除表中项目外，还检查了5分钟网格、项目序位连续、空腹10时前完成、床旁7号机和输出聚合标记。所有成功事件的项目数均等于输入服务项目数；所有未成功事件在结果表中保留失败原因“截止前无完整可行方案”。因此63.95\%不是项目级完成比例，也不是删除未排患者后的条件比例。

\subsection{补充敏感性情景}

赛题明确要求考虑峰/平配置和约5\%复杂病例；二者已进入主模型。为审核数据证据与假设边界，本文另设置七个\textbf{补充检验}，它们不是赛题指定情景：取消5\%时长上浮、医生容量采用训练日下四分位、只允许严格项目级设备证据、憋尿准备45/90分钟、跨室转运5/15分钟。每次只改变一个因素，结果见表\ref{tab:sensitivity}和图\ref{fig:sensitivity}。

\begin{table}[H]
  \centering
  \caption{主推荐策略的补充敏感性检验}
  \label{tab:sensitivity}
  \begin{tabular}{lrrrr}
\toprule
情景 & 住院48h率 & 背景预约日率 & 条件等待P90/h & 排入分钟 \\
\midrule
主情景：5\%上界时长 & 63.95\% & 59.70\% & 40.42 & 1,819,800 \\
无特殊时长上浮 & 63.97\% & 59.70\% & 40.42 & 1,819,660 \\
医生容量下四分位 & 61.96\% & 58.95\% & 42.00 & 1,769,275 \\
仅严格设备兼容 & 71.57\% & 49.67\% & 45.33 & 1,664,965 \\
憋尿准备45分钟 & 63.89\% & 59.88\% & 40.25 & 1,821,400 \\
憋尿准备90分钟 & 64.12\% & 59.17\% & 40.17 & 1,815,485 \\
跨室转运5分钟 & 63.93\% & 59.71\% & 40.50 & 1,821,800 \\
跨室转运15分钟 & 63.99\% & 59.76\% & 40.33 & 1,819,020 \\
\bottomrule
\end{tabular}

\end{table}

\begin{figure}[H]
  \centering
  \includegraphics[width=0.70\textwidth]{fig05_sensitivity.pdf}
  \caption{补充证据与参数边界相对主情景的变化}
  \label{fig:sensitivity}
\end{figure}

取消特殊时长上浮后住院率只增加0.013个百分点，背景率几乎不变。原因是多数项目上界与名义值在5分钟离散后差异有限，且系统瓶颈主要来自项目兼容和峰值到达。该结果不否定赛题的5\%条件，而是说明其在当前时长分档下边际影响较小。

医生容量降至下四分位后，住院率由63.95\%降至61.96\%，背景率由59.70\%降至58.95\%，住院首检等待中位数由4.50小时升至8.08小时，表明住院48小时目标对医生班次明显敏感。相比之下，单纯增加普通设备若没有医生同步到位，改善有限。

只允许严格设备证据时，背景率降至49.67\%，住院率反升至71.57\%。这不是模型变好，而是大量体检和门诊项目因缺少严格项目级匹配被直接判为不可排，释放了容量给住院。该现象说明项目能力覆盖必须用两类指标联合解释；部署前应由超声科逐条确认50个回退项目的可用设备，确认后再冻结正式能力矩阵。

憋尿准备45分钟时住院率为63.89\%，90分钟时为64.12\%，背景率分别为59.88\%和59.17\%。准备更晚会使部分背景事件无法在计划日完成，却为住院释放少量容量；准备更早又会改变排序与碎片。因此准备时长不是单调优化旋钮，应由预约通知、到院时间和实际准备状态共同决定。

跨室转运5、10和15分钟时，住院率分别为63.93\%、63.95\%和63.99\%，背景率分别为59.71\%、59.70\%和59.76\%。完成率在三点上变化较小但非单调，这是离散时隙、候选顺序和局部容量碎片共同作用的结果，不表示更慢转运能提升临床效率。上线时应用实测诊室动线替代该参数。

\subsection{结构合理性分析}

本文结果与诊断资源调度文献中的两个结论一致：一是混合患者资源应把事前容量保护与动态优先级结合\cite{green2006,patrick2007}；二是超声预约应同时考虑检查室分配和随机时长，而不是只给患者时间不给房间\cite{chen2022}。不同之处在于，本文没有用仿真分布代替题目缺失字段，而是以题给区间构造可解释上下界，并把项目能力证据强弱直接放入敏感性。

\section{模型的评价}

\subsection{模型优点}

\begin{enumerate}
  \item \textbf{口径完整。} 事件、项目、设备和医生均有明确单位；未分配患者保留在分母，避免把项目完成率冒充患者完成率。
  \item \textbf{约束统一。} 问题二和问题三共用同一房间、医生、患者、兼容、准备和转运规则，主模型与回放模型不存在可行域差异。
  \item \textbf{证据可追溯。} 设备严格匹配、床旁明确和类别回退分别记录；训练期与留出期物理隔离，未来报告时刻和历史派单不进入决策。
  \item \textbf{方案可执行。} 输出下沉到每位住院患者的项目次序、开始时刻和机房，并给出代表日完整甘特图，而非仅报告总容量或理论上界。
  \item \textbf{权衡透明。} 联合策略同时报告两类服务率和等待分位数；严格兼容情景的住院虚增被识别为服务覆盖下降，而不是误写为性能提升。
\end{enumerate}

\subsection{模型局限}

第一，历史数据缺少真实预约时刻、扫查起止和报告提交时刻。本文用题给区间确定绝对时长，用最后项目结束代理报告完成，因此无法分离扫查与报告环节。第二，未来医生只有总并发容量，没有姓名、项目资质和班表；模型保证总人数不超限，但不能保证某一时刻恰有具备专项资质的医生。第三，门诊和体检只构造回顾性计划日压力，不能还原真实预约提前期、迟到、爽约和插队。第四，题面没有给出背景服务下限，本文只能以共享FCFS服务率构造不劣基线；若医院给出更高阈值，策略选择会沿帕累托前沿改变。第五，构造算法在给定策略族中求最早可行方案，未证明对全年大规模整数模型达到全局最优；本文的“最佳”只指三种已比较策略中住院主指标最高且背景不低于基线。

\section{模型的改进与推广}

\subsection{数据与求解方法的改进}

若医院能够补充扫查开始、扫查结束、报告提交、预约时刻和医生项目资质，模型可作三项升级：其一，分别估计扫查时长和报告时长，以双阶段队列替代完成代理；其二，以医生—项目资质矩阵替代总并发容量，形成设备与医生双匹配；其三，用真实预约、迟到与爽约分布进行滚动随机优化，并把已确认预约冻结、未确认预约柔性调整。运行层面的刷新频率和冻结窗口应由叫号提前量、患者到院状态及系统响应时间标定，不在缺少现场数据时预设固定分钟数。

\subsection{医院内部推广与跨机构迁移}

模型的可推广部分是“事件—项目—双资源—时间窗”结构，而不是A医院的具体房号、日班时间或10分钟转运。在A医院内部，心电、CT、MRI等多资源检查可沿用患者完整性、能力兼容和截止期约束，但必须重新定义服务时长、准备要求和人员资质。迁移到其他医院时，应先以当地设备清单、班表和患者动线替换配置参数，再用本地时间顺序留出期评价；不应直接移植本文的63.95\%结果。

部署可采用“计划层—执行层—反馈层”闭环。计划层每日生成未来48小时容量保护；执行层读取到院、准备完成、急诊插入和设备停机状态；反馈层按日监控两类完成率、未排原因、等待分位数和改号次数。只有在前瞻试点中同时保持门诊/体检服务和住院48小时指标，才将新策略固化为正式规则。

\subsection{主要结论}

本文完成了超声预约问题的患者级联合建模。问题一得到训练期科室需求、11类项目时长、22台设备能力证据和星期—5分钟医生容量；床旁7号机是最明确的单点瓶颈。问题二在完整共享资源约束下输出住院患者项目顺序、时段和机房；问题三把门诊、体检与住院置于同一连续年度日历。以共享FCFS背景率为不劣基线时，资源可行FCFS在留出期取得住院63.95\%、背景59.70\%的完成率，是三个候选策略中住院主指标最高者；联合稀缺度以住院率降至63.31\%为代价把背景率提高到61.61\%，构成另一帕累托点。独立核验未发现约束违规，但医生容量和设备项目证据对结果影响明显。因而实施路线不是单纯增加设备，而是先确认正式背景服务阈值和项目能力清单、保护床旁与专项设备、保障高峰医生并发，再沿可接受的服务权衡选择患者级滚动排程策略。
